% M7 CONSTRAINTS PAPER --- the programme in its evaluator role.
% A perspective/methods manuscript: the boundary theorem + the classification
% theorem (M5) + the hosting theorem (M6) + the no-gos, packaged as an explicit
% CHECKLIST that any proposed model of emergent physics on a causal/discrete
% order must satisfy.  Claim discipline (charter): every constraint is about the
% CLASS, stated with its epistemic grade; the paper turns necessity theorems
% into a decision procedure for evaluating models --- the programme stops
% building models and starts evaluating them.  Compiles with pdflatex+bibtex
% (revtex4-2).  DRAFT.
\documentclass[aps,prd,reprint,nofootinbib,superscriptaddress,floatfix]{revtex4-2}
\usepackage{amsmath,amssymb,amsthm}
\usepackage{bm}
\usepackage{booktabs}
\usepackage{hyperref}

\newtheorem{constraint}{C}

\begin{document}

\title{What any causal-substrate model must satisfy:\\
structural constraints as a checklist for emergent physics}

\author{Miqueias Alves Mendes}
\email{mikalvesm@gmail.com}
\affiliation{Independent Researcher, Ibiapina, Cear\'a, Brazil}
\date{\today}

\begin{abstract}
A research programme that couples matter to a Lorentz-invariant causal order has,
over its life, accumulated a set of \emph{necessity theorems}: statements of the form
``no model of this class can do $X$.'' We argue these theorems are more useful read
forwards, as a checklist. Any proposal that builds physics on a causal or discrete
order---an emergent gauge field, an emergent scale, a discrete quantum-gravity
kinematics---must pass a fixed list of structural constraints, each a theorem with a
named joint where it can be broken. We state eight: a valence constraint (invariant rules
on a Poisson order have infinite expected coordination); a classification constraint
(finite-valence, two-dimensional, loop-carrying substrates exist \emph{only} off the
invariant class); a hosting constraint (the matter-sector form is an a priori function
of the compact target, fixed by its bi-invariant metric); a parity constraint (no
$P$-odd pure-geometry observable, even spontaneously); a scale constraint (the sprinkling
is the unique thinning fixed point, with no relevant direction); an arrow constraint
(the causal axis emerges, its orientation does not); a topology constraint (the protected
charges are the target's matter textures, no gauge defect); and a discrete-symmetry
constraint (a symmetry breaks iff it has an intrinsic carrier). We show how to apply the checklist by
evaluating four candidate substrates---a foliated triangulation, a rigid crystal, an
invariant emergent-photon claim, and a causal-set matter sector---reading off which
constraint each meets or violates without a single simulation. The role change is the
point: the same theorems that told one programme what \emph{it} could build, tell the
field what \emph{any} such model must satisfy. The four constraints that are the programme's own
classification results are proved here, with their computational certificates, so the checklist
stands on its own.
\end{abstract}

\maketitle

\section{The inverted use of a no-go}
\label{sec:intro}

Exclusion results are usually presented defensively---``we tried $X$ and proved it cannot
happen here.'' But a sharp no-go has a second, more valuable life. Once ``no
Lorentz-invariant causal order supports an emergent photon'' is a theorem rather than a
failed campaign, it stops being a fact about \emph{one} programme and becomes a filter on
\emph{every} proposal in a class: a new model that claims an emergent $U(1)$ on an
invariant discrete order must now say which lemma it breaks, or it is already refuted.

This paper collects the necessity theorems accumulated by a matter-on-causal-set
programme~\cite{CorePaper,WenComplement,MatterGravityPRD,SU3PRD}, states them as a
\emph{checklist}, and---for the four that are the programme's own classification results
(Constraints~\ref{c:sna},~\ref{c:hosting},~\ref{c:topology},~\ref{c:carrier})---proves them
here, so the checklist is self-contained: their assembly and computational certificates are
given in Appendix~\ref{app:proofs}. The shift is deliberate: the programme stops building
models and starts evaluating them. Each constraint below is (i) a theorem about the intrinsic
class---a measurable functional of a causal order $(\mathcal C,\prec)$ and a compact internal
field, no embedding coordinate---carrying its epistemic grade; (ii) equipped with a named
\emph{joint}, the single hypothesis a counterexample must break; and (iii) applicable to a
candidate model by inspection, without simulation. Three disciplines from the parent programme
carry over: statements are about the class, never about Nature; grades are explicit
([theorem], [sketch], [measured]); and the places the constraints do \emph{not} reach are
declared, not hidden (Sec.~\ref{sec:scope}).

\section{The eight constraints}
\label{sec:constraints}

Throughout, an \emph{intrinsic} rule is a measurable functional of $(\mathcal C,\prec,n)$
alone, with $n\colon\mathcal C\to X$ a field into a compact target $X=G/H$; the substrate
is a Poisson sprinkling of $\mathbb M^d$ unless a model declares otherwise.

\begin{constraint}[valence]
\label{c:valence}
{\rm[theorem]} Every non-trivial invariant pairwise rule on a Poisson-sprinkled causal
order has \emph{infinite} expected coordination. The invariant of a causal pair is the
proper time $\Delta\tau$; the locus of neighbours at fixed $\Delta\tau$ is the hyperboloid
$H^{d-1}=\mathrm{SO}^+(d-1,1)/\mathrm{SO}(d-1)$, of infinite volume, so by Campbell--Mecke
$\langle z\rangle=\rho\,\mathrm{Vol}(H^{d-1})\int\Delta\tau^{d-1}q(\Delta\tau)\,d\Delta\tau
=\infty$ for every rule $q\not\equiv0$~\cite{CorePaper}. The closure to non-pairwise rules
holds through the Palm calculus, up to the covering relation.
\end{constraint}

\noindent\textbf{Joint.} A finite-valence \emph{and} invariant model must break
Constraint~\ref{c:valence}: either the rule is not a function of the intrinsic pair
invariant, or the measure is not the Poisson (Lorentz-invariant) one. ``Local'' and
``invariant'' cannot both hold for an intrinsic rule on the sprinkling.

\begin{constraint}[classification]
\label{c:sna}
{\rm[theorem for the named classes; App.~\ref{app:sna}]} Call a substrate
\emph{string-net-admissible} (SNA) if its Hasse diagram has (1) finite valence, (2)
two-dimensional amenable (polynomial) ball growth, and (3) a percolating, finite-dimensional
space of loops---the three requirements an emergent gauge field needs~\cite{WenComplement}.
Then \emph{SNA holds only off the invariant class}: no Lorentz-invariant nor
covariant-random measure is SNA, while deterministic crystals and graded foliations can be.
The exclusion assembles four proven pieces: Poisson invariance forces infinite valence
(Constraint~\ref{c:valence}); exchangeable orders have no locally finite Hasse diagram;
covariant sequential growth is either hyperbolic (its sparse Hasse is locally a
Galton--Watson tree) or one-dimensional (posts cut it into finite blocks); and a
square-generated percolating lattice is graded, hence foliated, hence not boost-invariant.
The realizability side is explicit: a five-generator $\mathbb Z^2$ crystal and the graded
square lattice are SNA, and both are non-invariant.
\end{constraint}

\noindent\textbf{Joint.} An emergent gauge field on a discrete order requires an SNA
substrate; by Constraint~\ref{c:sna} that substrate is non-invariant. A model must
therefore \emph{declare its non-invariance}---a preferred foliation or a fixed
crystal---or exhibit an invariant SNA substrate, which would refute the classification.
The constraint has now crossed to the quantum side [measured]: in the analytically
continued Benincasa--Dowker path integral over 2D orders, \emph{no region of the phase
diagram is SNA}---the random phase fails finite valence logarithmically (the box-order
mechanism), the crystalline phase fails it linearly in $N$ (dense Kleitman--Rothschild-like
layers), and the first-order transition is a bimodal mixture of the two. Both phases are
loop-rich; both pay in valence: the classical dichotomy (finite valence \emph{or}
percolating loops, never both) reappears intact under the quantum weight. The tension
belongs to the order, not to the measure.

\begin{constraint}[hosting]
\label{c:hosting}
{\rm[theorem (assembled) + measured; App.~\ref{app:hosting}]} If the invariant coupling depends on $G$ only through
its bi-invariant (Killing) metric---as $E=-J\sum(1/\dim_f)\,\mathrm{Re}\,\mathrm{Tr}
(U_iU_j^\dagger)$ does---then the substrate cannot select among compact groups, and the
emergent \emph{form} is an a priori function of the target $X=G/H$: the Goldstone count is
$\dim X$, the topological charge lies in $\pi_3(G)=\mathbb Z$ (Bott), the multiplets are
irreducible representations of the unbroken $H$ (Peter--Weyl), and strong-coupling
confinement is present and \emph{center-independent}. Measured across $SU(3),SU(4),G_2$,
and---decisively for the coset structure---across $O(3)=SU(2)/U(1)=S^2$, whose Goldstone
count is $2=\dim S^2$, not $3=\dim SU(2)$~\cite{SU3PRD,MatterGravityPRD}.
\end{constraint}

\noindent\textbf{Joint.} A model that predicts a matter-sector count differing from the a
priori values of its declared target $X$ is either adding structure the bi-invariant coupling
cannot see---and must name the extra input---or is wrong. The group is input; the form is not.

\begin{constraint}[parity]
\label{c:parity}
{\rm[theorem]} No $P$-odd pure-geometry observable exists in the intrinsic class, and none
can arise spontaneously. Every time-orientation-preserving isometry, spatial reflections
included, induces the identity on isomorphism classes of $(\mathcal C,\prec,n)$, so a
$P$-odd functional $O$ satisfies $O=-O\equiv0$; spontaneous breaking is excluded because $P$
acts trivially on the whole configuration space, and a redundancy cannot break (an Elitzur
mechanism)~\cite{CorePaper}.
\end{constraint}

\noindent\textbf{Joint.} A model with emergent chirality of geometric origin must break the
intrinsic-action lemma---exhibit an intrinsic observable on which spatial reflection acts
non-trivially. The exact, never-restored chirality of the neutrino is structural evidence
against any causal-order origin of chirality.

\begin{constraint}[scale]
\label{c:scale}
{\rm[theorem]} The Poisson law is the exact fixed point and, at the level of laws, the global
attractor of thinning (random deletion), and it has no invariant relevant direction: homogeneous
truncated correlations $|x|^{-\alpha}$ are eigenfunctions with eigenvalue $p^{\alpha/d}<1$ for
every $\alpha>0$; the density $\rho$ is the single marginal direction, the theory's one
unit~\cite{CorePaper}. Emergence is what survives deletion (topology, forms); non-emergence is
what deletion dilates away (scales, correlations, loops).
\end{constraint}

\noindent\textbf{Joint.} A model claiming an emergent length or mass scale of substrate origin
must exhibit a \emph{relevant} direction under thinning---a correlation that grows rather than
dilates. Absent that, every internal ratio is locked and exactly one number is imported.

\begin{constraint}[arrow]
\label{c:arrow}
{\rm[theorem for the axis; measured for the arrow]} The comparability \emph{axis} emerges
(order duality $\prec\leftrightarrow\succ$ is the only $T$ action, and $T$-odd expectations
vanish in equilibrium), but the \emph{arrow}---a condensed direction of time---does not arise
spontaneously; it is an input. Unlike parity, $T$ has a carrier (order duality is intrinsic),
so the closure here is empirical, not kinematic, and is declared as such.
\end{constraint}

\noindent\textbf{Joint.} A model deriving a spontaneous arrow of time must exhibit the order
parameter and its condensation; the class provides the axis for free and charges for the arrow.

\begin{constraint}[topology]
\label{c:topology}
{\rm[theorem + certified; App.~\ref{app:topology}]} The substrate protects exactly the \emph{matter textures}
classified by the nontrivial homotopy of the compact target---$\pi_n(X)\neq0$ gives a
discrete, conserved charge, invariant under thinning (topology does not flow under deletion,
a corollary of Constraint~\ref{c:scale})---and forbids every \emph{gauge defect} (monopole,
gauge string), which needs an emergent gauge field ruled out by Constraint~\ref{c:sna}. The
monopole is doubly forbidden: $\pi_2(G)=0$ and no emergent $U(1)$. Certified along the ladder
$\pi_1$ (global vortex), $\pi_2$ (baby Skyrmion, charge invariant under deletion of $40\%$ of
sites), $\pi_3$ (Skyrmion/baryon).
\end{constraint}

\noindent\textbf{Joint.} A model claiming a magnetic monopole or a topological charge in a
sector with $\pi_n(X)=0$ must supply an emergent gauge field (against Constraint~\ref{c:sna})
or add homotopy the target does not have. The protected sectors are the target's, and no more.

\begin{constraint}[discrete symmetry]
\label{c:carrier}
{\rm[theorem; App.~\ref{app:carrier}]} A discrete spacetime symmetry is breakable---explicitly or
spontaneously---if and only if it has an intrinsic \emph{carrier}, an observable of
$(\mathcal C,\prec,n)$ on which it acts nontrivially. Parity acts as the identity (spatial
reflection is an order isomorphism), so it has no carrier and cannot break, even spontaneously
(Constraint~\ref{c:parity} is the $P$ case). Time reversal acts as order duality
$\prec\leftrightarrow\succ$---a genuine carrier---so it is not kinematically protected, and its
closure is empirical. Charge conjugation, an internal involution, has an internal carrier and may
break dynamically. The three chiral sectors of the Standard Model are precisely the three without
a geometric carrier.
\end{constraint}

\noindent\textbf{Joint.} A model breaking $P$ in pure geometry must exhibit a nonzero intrinsic
$P$-odd functional---i.e.\ a carrier reflection acts on. None exists in the class; the criterion
turns each discrete-symmetry claim into a search for its carrier.

\section{Applying the checklist}
\label{sec:apply}

The constraints are meant to be run against a proposal by inspection. Table~\ref{tab:eval}
evaluates four substrates. The verdicts are read off the constraints, not simulated.

\begin{table*}[t]
\caption{\label{tab:eval}Four candidate substrates against the checklist. Each verdict names
the joint the model must and does open; no simulation enters---every entry is a corollary of
Sec.~\ref{sec:constraints}.}
\begin{ruledtabular}
\begin{tabular}{p{3.4cm}p{13.2cm}}
Candidate & Verdict \\
\colrule
Foliated triangulation (CDT-like) & SNA \emph{possible}, but breaks C\ref{c:sna} by declaring a
preferred foliation---non-invariant by construction; re-enters lattice universality rather than
extending it. \\
Rigid causal crystal & SNA \emph{possible} (finite valence, percolating plaquettes), but breaks
C\ref{c:sna} as a deterministic crystal; excluded from the invariant class by
Bombelli--Henson--Sorkin. \\
Invariant emergent-photon claim & Refuted unless it breaks C\ref{c:valence} \emph{or}
C\ref{c:sna}: an invariant SNA substrate does not exist in the named classes; the proposal must
exhibit one or name the broken lemma. \\
Causal-set matter sector (this programme) & Passes all eight; the matter form is fixed a priori
by C\ref{c:hosting} once a compact target is chosen; existence (ordering at finite coupling) is
measured, not claimed to be derived. \\
\end{tabular}
\end{ruledtabular}
\end{table*}

Two lessons generalize. First, most ``new'' emergent-gauge proposals on discrete orders are
already located by C\ref{c:valence}--\ref{c:sna}: they either quietly break invariance (and
then live on the condensed-matter side of the signature boundary~\cite{WenComplement}) or they
claim the impossible intersection and owe a broken lemma. Second, a model that survives
C\ref{c:valence}--\ref{c:scale} still imports exactly what the constraints say it must---one
density unit (C\ref{c:scale}), a compact target (C\ref{c:hosting}), an arrow (C\ref{c:arrow})---so
the checklist doubles as an inventory of the \emph{necessary} external inputs of any such model.

\section{Scope: where the checklist does not reach}
\label{sec:scope}

The constraints are theorems about a class, not about Nature, and their boundaries are part of
the statement. Constraint~\ref{c:sna} is proved for the named invariance classes
(Poisson, exchangeable, sequential-growth, embedded) with the crystalline realizability explicit;
genuinely non-Markovian dynamical geometry---a fluctuating causal measure with memory---is outside
it and is the one live frontier. Constraint~\ref{c:hosting} fixes the \emph{form} of the matter
sector, not its \emph{existence} (ordering is dynamical and measured), and assumes a non-degenerate
bi-invariant metric. Constraint~\ref{c:arrow} is empirical on the arrow. The checklist evaluates
kinematic admissibility; it does not predict phase diagrams, and a model can pass every constraint
and still be dynamically empty. An exclusion checklist that rejected every model would be as
useless as one that admitted every model; this one does neither, and it says exactly where each
verdict comes from.

\section{Outlook: from a theory of a class to a theory of constraints}
\label{sec:outlook}

Eight constraints are a beginning, not a closure. Constraints~\ref{c:topology}--\ref{c:carrier}
are already of this second generation---a classification of the protected topological sectors
(matter textures $\pi_n(X)\neq0$ survive coarse-graining; gauge defects are forbidden) and a
classification of the permitted spacetime-symmetry breakings by the carrier criterion ($P$
never, $T$ empirically, $C$ dynamically)---and the same method generates more: a classification
of the hostable compact groups (the form is fixed a priori by the bi-invariant metric), and, at
the live frontier, the admissibility of genuinely non-Markovian dynamical geometry, the one class
the constraints do not yet reach. Each new necessity theorem adds a row to the checklist.
The programme's centre of gravity thereby moves from asking \emph{what emerges from this substrate}
to asking \emph{what any substrate must satisfy to emerge a given physics}---the inverse problem,
posed rigorously as necessity rather than open-ended synthesis. The distinctive contribution is not
the ambition, which is old, but the method: constraints that are proved, graded, and attackable at
named joints. A field of emergent-spacetime proposals that adopted such a checklist would spend less
effort simulating models that a one-line corollary already excludes.

\begin{acknowledgments}
The parent measurements are reported, with full pre-registrations, death criteria and negative
results, in the programme's public record~\cite{CorePaper,WenComplement,MatterGravityPRD,SU3PRD}.
This manuscript reframes them as a checklist and proves the four classification constraints in
Appendix~\ref{app:proofs}; it introduces no new measurement.
\end{acknowledgments}

\appendix

\section{Proofs and certificates of the classification constraints}
\label{app:proofs}

Constraints~\ref{c:valence},~\ref{c:parity},~\ref{c:scale},~\ref{c:arrow} are proved in the
core paper~\cite{CorePaper}. Here we give the assembly and the computational certificates for
the four constraints that are the programme's own classification theorems. Each certificate is
a deterministic gate over the intrinsic data; grades are stated per step.

\subsection{Classification (Constraint~\ref{c:sna}): SNA $\Leftrightarrow$ non-invariant}
\label{app:sna}

\emph{Exclusion} [theorem, assembled]. Let a substrate be SNA (finite valence; amenable
two-dimensional ball growth; percolating finite-dimensional loops). Each named invariance class
fails one clause. (i) A Lorentz-invariant Poisson order has $\langle z\rangle=\infty$ by
Constraint~\ref{c:valence}---clause 1 fails. (ii) An exchangeable random order has a
conditionally constant covering density (Aldous--Hoover/Janson), so its Hasse diagram is either
a.s.\ empty or of infinite expected degree---clause 1 fails. (iii) Covariant sequential growth
splits: its \emph{sparse} regime ($p=\lambda/N$) has, by the covering lemma below, a Hasse
diagram that is a subgraph of the generating Erd\H{o}s--R\'enyi graph $G(N,\lambda/N)$, whose
local weak limit is a Galton--Watson tree, hence hyperbolic (exponential ball growth)---clause 2
fails; its \emph{dense} regime has posts of positive density, and no Hasse edge crosses a post,
so the order is an ordinal sum of finite blocks---one-dimensional, clause 2 fails. (iv) A
square-generated percolating lattice has balanced $4$-cycles, so the circulation
$\varphi(\gamma)=(\text{up}-\text{down})$ annihilates the cycle space; the potential
$r(x)=\varphi(x_0\!\to\!x)$ is well defined, grading the order into an intrinsic foliation, which
breaks boost invariance.

\emph{Covering lemma} [theorem]. In transitive percolation, every Hasse (covering) edge is a
\emph{directly sampled} edge of the generating graph: a covering pair $i\lessdot k$ has no
intermediary, so it cannot arise from a directed path of length $\ge2$; it is a length-one edge.
Hence $\mathrm{Hasse}\subseteq G(N,p)$, and a subgraph of a locally tree-like graph is locally
tree-like. Certified: over four regimes $\times$ six seeds, every Hasse edge lies in the
generating graph (zero violations); the sparse four-cycle density falls as $\sim1/N$
($0.072\!\to\!0.011$) against a constant $\approx1.9$ for a genuine two-dimensional lattice.

\emph{Realizability} [construction]. A $\mathbb Z^2$ crystal with five generators, and the graded
square lattice, are SNA (finite valence, percolating plaquettes, amenable growth) and manifestly
non-invariant (a deterministic crystal is not Lorentz-invariant by Bombelli--Henson--Sorkin; the
graded lattice is foliated). \emph{Certificate.} A unified battery placing one representative of
each class on the same measurement axis returns SNA true \emph{only} in the two non-invariant
rows (crystal, foliated lattice), with no invariant or covariant counterexample.

\subsection{Hosting (Constraint~\ref{c:hosting}): the form is a priori in $X=G/H$}
\label{app:hosting}

\emph{Mechanism} [theorem]. The coupling $E=-J\sum(1/\dim_f)\mathrm{Re}\,\mathrm{Tr}
(U_iU_j^\dagger)=-J\sum v_i\!\cdot\!v_j$ is a functional of the bi-invariant (Killing) inner
product on $G$, which every compact group carries. The substrate therefore cannot distinguish
groups beyond dimension, coset structure, and topology, and the emergent form is fixed a priori:
the Goldstone count is $\dim(G/H)$ (Goldstone's theorem), the topological charge lies in
$\pi_3(G)=\mathbb Z$ (Bott, for every simple compact $G$), the unbroken multiplets are
irreducible representations of $H$ (Peter--Weyl), and strong-coupling confinement is present and
center-independent (the character expansion does not see $Z(G)$).

\emph{Certificate.} The Goldstone count equals $\dim X$ across $SU(2)\!\to\!3$,
$SU(3)\!\to\!8$, $SU(4)\!\to\!15$, $G_2\!\to\!14$, and---the discriminating case---the coset
$O(3)=SU(2)/U(1)=S^2\!\to\!2$, which is $\dim S^2$ and \emph{not} $\dim SU(2)=3$: the form tracks
the coset $X$, not the group $G$. The $\pi_3$ charge ladders of $SU(2)$ and $SU(3)$ coincide
digit for digit ($0.795,0.892,0.951\to1$; anti-charge $=-B$), confirming $\pi_3=\mathbb Z$
independent of $N$. Untested groups follow a priori ($SU(5)\!:\!24$, $F_4\!:\!52$, $E_6\!:\!78$;
$\pi_3=\mathbb Z$).

\subsection{Topology (Constraint~\ref{c:topology}): protected textures, forbidden defects}
\label{app:topology}

\emph{Statement} [theorem + certificate]. A matter texture of degree $n$ (a charge of the field
$n\colon\mathcal C\to X$) is protected iff $\pi_n(X)\neq0$; the charge is a discrete invariant,
conserved and---being topological---unchanged under thinning (Constraint~\ref{c:scale}: only
non-topological data flows). A gauge defect (monopole, gauge string) needs an emergent gauge
field, forbidden by Constraint~\ref{c:sna}; the monopole is doubly forbidden ($\pi_2(G)=0$ and no
emergent $U(1)$). \emph{Certificate.} Integer charge along the homotopy ladder: a $\pi_1$ global
vortex (winding $+1,+2,-1$), a $\pi_2$ baby Skyrmion (Berg--L\"uscher charge $+1,+2,-1$), and a
$\pi_3$ baryon (convergent ladder, anti-charge $=-B$). The $\pi_2$ charge is invariant under
random deletion of up to $40\%$ of the sites ($Q=0.9996$ down to keep-fraction $0.6$)---topology
does not flow.

\subsection{Discrete symmetry (Constraint~\ref{c:carrier}): breakable iff carried}
\label{app:carrier}

\emph{Statement} [theorem]. A discrete symmetry $S$ acting on $(\mathcal C,\prec,n)$ is breakable,
explicitly or spontaneously, iff there is a nonzero intrinsic $S$-odd functional (a carrier). If
$S$ acts trivially (an automorphism of the configuration), every $S$-odd $O$ obeys $O=-O\equiv0$,
and spontaneous breaking is impossible for want of an order parameter (an Elitzur argument).
\emph{Application.} Spatial reflection preserves $\prec$ (it is an order isomorphism), so parity
acts as the identity---no carrier, closed both channels, kinematically. Time reversal acts as
order duality $\prec\leftrightarrow\succ$, a genuine carrier, so it is not kinematically
protected; the closure is empirical. Charge conjugation is an internal involution with an internal
carrier and may break dynamically. \emph{Certificate.} Over twenty sprinklings, reflection
preserves the order-isomorphism class $20/20$ and the mean of a candidate $P$-odd functional is
compatible with zero; the intrinsic $T$-odd functional $\sum_i(\mathrm{outdeg}_i^2-
\mathrm{indeg}_i^2)$ is nonzero per causet and flips sign under dualization $20/20$ (a carrier),
while its mean is compatible with zero (the law is self-dual); the $C$-odd functional flips sign
under the internal involution. The three chiral sectors of the Standard Model are the three
lacking a geometric carrier.

\begin{thebibliography}{9}
\bibitem{CorePaper} M. Alves Mendes, \emph{What a Lorentz-invariant discrete order can carry: an
axiomatic core for emergence on causal substrates}, companion paper (2026).
\bibitem{WenComplement} M. Alves Mendes, \emph{Why Lorentz-invariant causal sets cannot support an
emergent $U(1)$: a causal complement to string-net theory}, companion paper (2026).
\bibitem{MatterGravityPRD} M. Alves Mendes, \emph{Topological matter and emergent gravitation from a
causal network}, companion paper (2026).
\bibitem{SU3PRD} M. Alves Mendes, \emph{Colour ferromagnetism in a causal network: confinement, the
meson octet, and the order of the transition}, companion paper (2026).
\end{thebibliography}

\end{document}
